\chapter{Computing the Posterior}\label{ch:computation}

\begin{tcolorbox}[colback=gray!6,colframe=gray!50,title=In brief,fonttitle=\normalfont\itshape,coltitle=black,colbacktitle=gray!15,boxrule=0.4pt,arc=1pt]
Horseshoe Gibbs samplers took under a second in the two Gaussian outcome designs with 51 and 56 columns, half the time of a cross fitted lasso in one and a twentieth of the time of a cross fitted random forest in the other, so for models of that size computation is no reason to prefer double machine learning\index{double machine learning!Bayesian alternative to}. With 300 or 400 covariates the comparison turns, and DML with a lasso nuisance is the cheaper of the two (Section~\ref{comp:sec:engine}). When the likelihood cannot be evaluated cheaply, a quantile network trained on simulated pairs of parameter and data turns uniform noise and a summary of the data into posterior draws, and the cost moves to a single training phase. Diagnostics are part of the computation, and a closing table says which engine to use when.
\end{tcolorbox}

One question organizes this chapter: how long it takes to get a few thousand posterior draws that we can trust. The answer has two parts, a time and a check. A Gibbs sampler serves when the likelihood is cheap to evaluate; when it is not, generative Bayesian computation\index{generative Bayesian computation|textbf}\index{GBC|see{generative Bayesian computation}} learns a map from noise and a summary of the data.

% comp s01, from old m16 line 7
\section{Why computation is the objection}\label{comp:sec:why}

Cheap computation is one advantage of double machine learning\index{double machine learning}\index{double machine learning!computational cost|textbf}. A cross fitted DML estimate needs a handful of regressions; a posterior needs thousands of Markov chain iterations\index{computational cost of a posterior}. The horseshoe\index{horseshoe prior} sampler is fast enough for the models of this book (Section~\ref{comp:sec:cost}).

A posterior is a set of draws, and every target in this book is a function of the draws, such as a risk difference averaged over a population $Q$ that the analyst chooses (Chapter~\ref{ch:mccullagh}). A sampler is therefore indifferent to the target; a learned map is trained for the parameter we ask of it (Section~\ref{comp:sec:core}).



The quantile network plays two roles in this book: here it learns a posterior from simulated pairs of parameter and data; in Chapter~\ref{ch:gente} it regresses the outcome on covariates and treatment using observed data, with uncertainty from the Bayesian bootstrap. Chapter~\ref{ch:blp} uses a variant: a regression network learns a posterior mean from simulated pairs, and residuals from cross fitting supply the spread.

% comp s02, from old m16 (cost, data augmentation, approximations, sampler)
\section{What a posterior costs}\label{comp:sec:cost}

Write $Z$ for the design matrix and $\sigma^2\Lambda$ for the prior covariance of the coefficients, with $\Lambda=\mathrm{diag}(\lambda_j^2\gscale^2)$. Each iteration of the Gibbs sampler\index{Markov chain Monte Carlo|textbf}\index{MCMC|see{Markov chain Monte Carlo}} of Chapter~\ref{ch:horseshoe} draws the coefficients from a Gaussian full conditional, which needs one Cholesky factorization\index{Cholesky factorization|textbf} of the $p\times p$ matrix $A=Z^\top Z+\Lambda^{-1}$ at $O(p^3)$, and then draws the scales and the error variance from inverse gamma full conditionals at $O(p^2)$ at most. The cost of a posterior\index{computational cost of a posterior|textbf} is one factorization times the number of iterations: the first factor is known from the size of the model, the second has to be checked.

\subsection{The horseshoe Gibbs sampler}\label{sec:hsgibbs}

The data enter the sampler\index{Gibbs sampler|textbf}\index{Gibbs sampler!for the horseshoe|textbf}\index{horseshoe prior!Gibbs sampler|textbf} only through the cross products $Z^\top Z$, $Z^\top Y$ and $Y^\top Y$, which are formed once at $O(np^2)$. With $A = LL^\top$, a draw of the coefficients is $A^{-1}Z^\top Y + \sigma L^{-\top}\eta$ with $\eta \sim N(0, I_p)$, and the residual sum of squares that the update of $\sigma^2$ needs is $Y^\top Y - 2\beta^\top Z^\top Y + \beta^\top Z^\top Z \beta$. Nothing in an iteration then touches $n$, and for a Gaussian outcome the sampler is as fast for ten thousand units as for two hundred. The simulation sampler timed below forms the residual vector $Y - Z\beta$ in each iteration instead, an $O(np)$ step that is small at $n = 200$.

The scales mix slowly when many coefficients are near zero, so the number of iterations has to be checked and not assumed. Most simulations of this book use 1,500 iterations with the first 500 discarded, the 401(k) analysis 3,000 with the first 1,000 discarded, and the horseshoe fits inside double machine learning in Section~\ref{sec:baipw} 600 sweeps with the first 200 discarded; the checks of Section~\ref{comp:sec:checks} say whether that was enough.

In the simulation of Section~\ref{sec:ricsim}, with $n = 200$ and $p = 51$ columns (the treatment and 50 covariates), 1,500 iterations took 0.31 seconds on one core of an Intel Xeon Gold 6240R at 2.4 GHz, the mean of 20 timed runs (the slowest took 0.33 seconds). One cross validated lasso fit took 0.064 seconds, so the sampler costs about five lasso fits, and the ten lasso fits of a five fold cross fitted\index{cross fitting} estimate took 0.60 seconds, twice the sampler. On the 401(k) data of Chapter~\ref{ch:k401}, with $n = 9{,}915$ and 56 columns (the treatment, an intercept and the 54 terms of a quadratic basis in the nine household characteristics), 3,000 iterations took 0.63 seconds (the mean of five runs, 0.58 to 0.64). On the same data one lasso fit took 0.09 seconds, one random forest fit 1.2 seconds, and the ten forest fits of one cross fitted estimate 12 seconds, about nineteen times the sampler. The lasso is cross validated with five folds, the forest has 200 trees, at least 5 observations per leaf and 4 covariates per split (the settings of the DML comparison in that chapter), and both are fitted on four fifths of the rows, as in one fold. Appendix~\ref{app:code} names the script. For large $p$ the factorization dominates, and when $p > n$ the exact sampler of Bhattacharya, Chakraborty and Mallick (2016) costs $O(n^2 p)$ through the Woodbury identity, orders of magnitude less for $p$ in the thousands and $n$ in the hundreds.

We checked the treatment coefficient, the quantity reported in the tables. In each of 20 data sets of the design of Section~\ref{sec:ricsim} we ran four chains of 1,500 iterations from the sampler's default start, discarded the first 500, and computed the split $\hat R$\index{split R hat@split $\hat R$|textbf} and the effective sample size\index{effective sample size|textbf} from the 4,000 draws left. For the independent horseshoe the largest $\hat R$ was 1.003 and the smallest effective sample size 1,660, and for the coupled horseshoe 1.007 and 1,340. In 10 data sets each of the designs with $p = 150$ and $p = 300$ the largest $\hat R$ was 1.012 and the smallest effective sample size 620, and only one of the 80 simulated data sets, the coupled prior at $p = 150$, had an $\hat R$ above 1.01. In 10 data sets of the design with thirty weak confounders, run the same way, both coupled priors passed in every data set, with the largest $\hat R$ 1.004 and the smallest effective sample size 769, and the independent horseshoe did not: six of the ten had an $\hat R$ above 1.01, the largest 1.020, and eight had fewer than 400 effective draws, the smallest 234. With 234 effective draws the Monte Carlo error of a posterior mean is about a fifteenth of a posterior standard deviation, far less than that design's bias of 0.670, which none of the independent horseshoe's intervals covers; the endpoints of those intervals are less secure. On the 401(k) data, four chains of 3,000 iterations with the first 1,000 discarded gave $\hat R$ of 1.001 for the independent prior and 1.000 for the coupled prior, and an effective sample size equal to the number of draws. Apart from the 600 sweep runs of Section~\ref{sec:baipw}, which we did not check, and the independent horseshoe with thirty weak confounders, the chains behind Chapters~\ref{ch:horseshoe} and \ref{ch:k401} mixed for the treatment coefficient; the check says nothing about the other parameters.

Coefficients that must stay unshrunk, such as an intercept, the treatment coefficient and an estimated propensity, enter $A$ with a prior precision of essentially zero, $10^{-8}$ in the 401(k) analysis of Chapter~\ref{ch:k401}. Unless the outcome and every column are centred, a design without the intercept forces the regression through the origin, and the treatment coefficient, which multiplies the only uncentred column, absorbs the outcome mean. Our first run made this mistake and produced an effect more than twice the correct value, with a deceptively narrow posterior.

\subsection{Data augmentation for binary and count outcomes}

The logistic likelihood is not conjugate to a Gaussian prior, but the Polya Gamma\index{Polya Gamma@P\'olya Gamma} augmentation\index{Polya Gamma@P\'olya Gamma!augmentation|textbf}\index{data augmentation|textbf} of Polson, Scott and Windle (2013) makes it conditionally Gaussian via the identity
\[
\frac{(e^{\zeta})^a}{(1 + e^{\zeta})^b} = 2^{-b}e^{\tilde y\zeta}\int_0^\infty e^{-\omega\zeta^2/2}\,p(\omega \mid b, 0)\,d\omega, \qquad \tilde y = a - b/2,
\]
where $p(\omega\mid b,0)$ is the Polya Gamma density (Appendix~\ref{app:toolkit}). A logistic unit with linear predictor $\zeta$ has $a=y_i$ and $b=1$, so $\tilde y_i=y_i-1/2$ as in Chapter~\ref{ch:horseshoe}. Given the $\omega_i$ the likelihood is a Gaussian regression with weights $\omega_i$ and working response $\tilde y_i/\omega_i$, so the horseshoe step\index{horseshoe prior!in logistic regression|textbf} applies unchanged, and given the linear predictor each $\omega_i$ has a Polya Gamma full conditional that can be sampled exactly. The same device handles negative binomial counts and multinomial outcomes.

Alternating the two draws samples the propensity model of Chapter~\ref{ch:horseshoe}, with the coupled priors that share scales between the treatment and outcome equations, and the coefficients of a binary outcome such as cancer in Chapter~\ref{ch:smoking}. Now $Z^\top\Omega Z+\Lambda^{-1}$, with $\Omega=\mathrm{diag}(\omega_i)$, depends on $\omega$, so $Z^\top\Omega Z$ is formed in every iteration at $O(np^2)$, and unlike the Gaussian sampler this one depends on $n$.

\subsection{When the sampler is too slow}

Two approximations are standard. The Laplace approximation\index{Laplace approximation|textbf} replaces the posterior by a normal centred at the mode with covariance from the Hessian; integrated nested Laplace approximations\index{integrated nested Laplace approximation|textbf}\index{INLA|see{integrated nested Laplace approximation}} (Rue, Martino and Chopin 2009) extend it to latent Gaussian models. Variational Bayes\index{variational inference|textbf} approximates the posterior by the closest member, in Kullback and Leibler divergence, of a tractable family.

For a treatment coefficient with a flat prior and a large sample, the Bernstein von Mises theorem makes the posterior close to normal under regularity conditions, and the Laplace approximation is then good. It is poor for the horseshoe scales, whose posteriors are skewed and bounded below, but those are nuisance parameters that are integrated out. Mean field factorization removes the posterior correlation that confounding adjustment creates between treatment and confounder coefficients, underestimating variance. An effect interval can be too short by a factor that depends on confounding strength, and nothing in the optimization reports it.

Either approximation can supply starting values or screen covariates before a full run, but the reported analyses use neither: every chain starts from the sampler's default values. A tail probability may still be wrong under an approximation that is good at the centre.

When the log posterior can be evaluated and differentiated but the full conditionals have no closed form, Hamiltonian Monte Carlo (HMC)\index{Hamiltonian Monte Carlo|textbf} is the general sampler. It adds a Gaussian momentum, follows the gradient by leapfrog steps and accepts or rejects the end point by a Metropolis step, so the normalizing constant is never needed; the no U turn sampler of Hoffman and Gelman (2014) chooses trajectory lengths during sampling and adapts the step size during warm up. Positive scales are sampled on the log scale with the Jacobian. For the horseshoe, the non centred form $\beta_j=\sigma\gscale\lambda_jz_j$ with $z_j\sim N(0,1)$ can ease funnel geometry, especially when coefficients are weakly informed by the data. Divergent transitions call for a check of posterior geometry; the checks of Section~\ref{comp:sec:checks} still apply. Each leapfrog step needs a gradient that touches all $n$ rows, so the cost is iterations times steps times $np$. 

\paragraph{Running the analyses at scale.} For many outcomes, subgroups or sites, batch the Cholesky factorizations\index{Cholesky factorization} across outcomes that share a design matrix, and start each run at the posterior mean of a related run, which shortens burn in. When one model is fitted to many similarly structured data sets, a generative map\index{generative Bayesian computation!cost|textbf} replaces each run by a forward pass (Section~\ref{comp:sec:algorithm}), with the calibration check of Section~\ref{comp:sec:checks} run once on the trained network.

% comp section 03 (comp:sec:noise), phase 2 rewrite; Knothe, Rosenblatt, Brenier and ordering advice kept from m14, generator display now the triangular map, amortization sentence in comp/s05; figure by book_scripts/comp_figs.py
\section{A posterior is a map from noise}\label{comp:sec:noise}

Every first course in simulation teaches one trick\index{inverse transform sampling|textbf}. To draw from a distribution with distribution function $F$, draw $U$ uniform on $(0,1)$ and return $F^{-1}(U)$: the result is at most $z$ exactly when $U\le F(z)$, which has probability $F(z)$. A posterior is no different, except that nobody hands us the inverse. Yet a function of this kind always exists. That is noise outsourcing\index{noise outsourcing|textbf}: one random object can be generated from another by feeding it, with an independent uniform number, to one fixed function.

\begin{theorem}[Noise outsourcing]\label{comp:thm:outsourcing}
Let $X$ and $Y$ be random elements of standard Borel spaces. On a possibly enlarged probability space there are a random variable $U$, uniform on $(0,1)$ and independent of $X$, and a measurable function $G$ such that $Y = G(U,X)$ almost surely\index{noise outsourcing}.
\end{theorem}

This is the transfer theorem\index{Kallenberg transfer theorem|see{noise outsourcing}}\index{transfer theorem|see{noise outsourcing}} of Kallenberg (2002). For real $Y$ one can take $G(u,x)=F^{-1}(u\mid x)$, the conditional quantile function\index{quantile function|textbf}, with $F^{-1}(u\mid x)=\inf\{y:F(y\mid x)\ge u\}$ the generalized inverse, so that atoms are covered: at a jump of $F$ the level $U$ is spread uniformly over the jump by an extra independent uniform variable, which is why the space may have to be enlarged (Exercise~\ref{ex:ch08:1}; the proof is in Appendix~\ref{app:theory}). Other standard Borel spaces reduce to the real case by a Borel isomorphism. With $Y=\theta$ and $X=S(D)$ for a data summary $S$, $G(U,s)$ has law $\pi(\theta\mid S=s)$ for almost every $s$, the full posterior when $S$ is sufficient (Section~\ref{comp:sec:summaries}).

The theorem is an existence statement. $G$ changes when the prior or the model does, and outside conjugate cases it has no formula. It is not unique either, since any measure preserving shuffle of the levels gives another function with the same output law, but only one choice increases in $u$, the quantile function\index{quantile function!posterior|textbf}, which for scalar $\theta$ is the one we learn. A posterior is then fully described by a single increasing function on $(0,1)$ (Figure~\ref{comp:fig:inverse}). Its slope is the reciprocal of the posterior density, shallowest at a mode and steep in the gap between two.

\begin{figure}[ht]
\centering
\includegraphics[width=\linewidth]{comp_inverse.pdf}
\caption{One bimodal posterior, the mixture $0.4\,N(-1.5,0.45^2)+0.6\,N(1.5,0.6^2)$. Left: its distribution function, with six uniform draws carried across to the curve and down to draws of $\theta$. Centre: the quantile function. Right: 5{,}000 draws $Q(U)$ and the exact density.}
\label{comp:fig:inverse}
\end{figure}

For a vector $\theta=(\theta_1,\ldots,\theta_d)$ we apply the scalar statement one coordinate at a time, with independent uniforms $U_1,\ldots,U_d$:
\[
\theta_1 = G_1(U_1,s),\qquad \theta_j = G_j(U_j,s,\theta_1,\ldots,\theta_{j-1}),\quad j=2,\ldots,d,
\]
where $G_j$ is the quantile function of $\theta_j$ given $s$ and the earlier coordinates. The map is triangular\index{triangular map|textbf}. It is the factorization $p(\theta_1\mid s)\,p(\theta_2\mid s,\theta_1)\cdots$ turned into a sampler, which transport theory calls the Knothe and Rosenblatt map\index{Knothe Rosenblatt map|textbf} (Rosenblatt 1952; Knothe 1957). Brenier's theorem (Brenier 1991) gives a different map, the gradient of a convex function; the two generally differ, coinciding when the coordinates are independent. For computation any map with the right output law will do, and the triangular one reduces a vector problem to $d$ scalar quantile regressions.

Every coordinate order gives an exact generator, but early errors pass to later coordinates. We put the quantity of interest first, learning the treatment effect from its marginal posterior given the summary before the nuisances. No ordering repairs bias from a prior that ignores the propensity score (Section~\ref{comp:sec:core}).

The theorem allows atoms. A spike and slab posterior puts positive probability on a coefficient being exactly zero, and its quantile function has a flat piece at zero whose length on the level axis is that probability. The horseshoe posterior has a density, so its quantile function has no flat piece.

Pushing several coordinates through their own marginal quantile functions with one shared $U$ gives the right margins and the wrong joint law: the coordinates become comonotone, perfectly rank correlated whatever the posterior says. Driving the treated and control quantile functions with one uniform is the same mistake in causal dress, and it assumes rank preservation\index{rank preservation} (Section~\ref{sec:ot}).

One map serves every question through an identity. As $G(U,s)$ with $U$ uniform on $(0,1)^d$ has the posterior law, for any $h$ with $\E\lvert h(\theta)\rvert<\infty$
\begin{equation}\label{comp:eq:quantileidentity}
\E\{h(\theta)\mid s\} = \int_{(0,1)^d} h\bigl(G(u,s)\bigr)\,du.
\end{equation}
For $d=1$, $h(\theta)=\theta$ gives the posterior mean as the area under the quantile function\index{quantile function}, and $h(\theta)=\one(\theta>c)$ gives the posterior probability of $\theta>c$ as the share of levels at which the quantile function exceeds $c$. The Wasserstein-1 distance between two scalar posteriors is the area between their quantile functions. A utility is another $h$. A risk criterion reweights the levels of the utility's own quantile function (Chapter~\ref{ch:cdt}): for scalar $\theta$ that is the utility of the quantile function of $\theta$ only when the utility increases in $\theta$, and in general the utilities of the draws are sorted first. Once $G$ is learned a new question costs an integral, not a new fit: the learned map is the model, and each integral is accounting (Chapter~\ref{ch:mccullagh}).

% comp section 04 (comp:sec:network, gbc:sec:qnn), phase 2 rewrite; "The proposition is exact" is in comp/s06, "A check that can be run" in comp/s09
\section{Learning the map: the check loss and the quantile network}\label{gbc:sec:qnn}\label{comp:sec:network}

We cannot write $G$ down, but regression on simulated pairs of parameter and summary learns a conditional mean under squared error and a conditional quantile under the check loss\index{check loss|textbf}; give the network the level as an input and it learns every quantile at once.

\subsection{The quantile neural network}

A quantile neural network\index{quantile neural network|textbf} is a function $f_\phi(u,s)$ of a level $u\in(0,1)$ and a value $s$ of the summary. We train it on simulations from the prior predictive distribution: draw $\theta^{(j)}\sim\pi$ and $D^{(j)}$ from the model, set $s^{(j)}=S(D^{(j)})$, draw a level $u_j$ uniformly, and minimize the check loss\index{quantile neural network!training by check loss|textbf}
\begin{equation}\label{gbc:eq:pinball}
\mathcal L(\phi) = \frac1J\sum_{j=1}^J \rho_{u_j}\bigl(\theta^{(j)} - f_\phi(u_j, s^{(j)})\bigr),\qquad \rho_u(r) = r\{u - \one(r<0)\}.
\end{equation}
At inference $f_\phi(V,s)$ with $V$ uniform and $s=S(D^{\mathrm{obs}})$ is an approximate posterior draw. For a vector $\theta$ the loss is applied coordinate by coordinate in the triangular order, with the earlier coordinates added to the inputs.

\begin{proposition}[What the check loss learns]\label{gbc:prop:pinball}
Let $(\theta, S)$ have the joint law induced by the prior and the model, with $\E\lvert\theta\rvert<\infty$, and let $U$ be uniform and independent of them. Among measurable functions $f(u,s)$ with $\E\lvert f(U,S)\rvert<\infty$, the expected check loss $\E\,\rho_U(\theta - f(U,S))$ is minimized by $f^\ast(u,s) = Q(u\mid s)$, the conditional quantile function of $\theta$ given $S=s$, and the minimizer is unique for almost every $(u,s)$ at which the conditional distribution function is strictly increasing.
\end{proposition}
\begin{proof}
Conditioning on $(U,S)$, it suffices to minimize $g(q) = \E[\rho_u(\theta - q)\mid S=s]$ over $q$. The function $g$ is convex with right derivative $F(q\mid s) - u$, which is nonnegative exactly when $q\ge Q(u\mid s)$, so the minimum is attained at the conditional $u$ quantile, uniquely where $F(\cdot\mid s)$ is strictly increasing. Integrating over $(u,s)$ gives the result.
\end{proof}

The proposition is about the population minimizer. For a trained network the evidence is the checks of Section~\ref{comp:sec:checks}. Integrated over $u$, the check loss of a curve that increases in $u$ is half the continuous ranked probability score\index{continuous ranked probability score|textbf}\index{CRPS|see{continuous ranked probability score}} of the law the curve describes, a strictly proper score (Gneiting and Raftery 2007). 

\subsection{Quantile neural networks explained}\label{gbc:sec:qnnexplained}

The check, or pinball\index{pinball loss|see{check loss}}, loss charges $u$ per unit when the target lies above the guess and $1 - u$ when below, so the guess that balances the two expected charges is the $u$ quantile (Figure~\ref{fig:qnnloss}). Classical quantile regression fits one level at a time (Koenker and Bassett 1978). The implicit quantile network of Dabney, Ostrovski, Silver and Munos (2018), introduced for distributional reinforcement learning, takes the level as an input, so one network\index{quantile neural network} represents the whole quantile function $u \mapsto Q(u \mid s)$.

\begin{figure}[ht]
\centering
\includegraphics[width=\linewidth]{qnn_loss.pdf}
\caption{Left: the check loss for levels 0.1, 0.5 and 0.9. Right: expected check loss for a Gamma(2) target; the minimizers, marked by vertical lines, are the 0.1, 0.5 and 0.9 quantiles.}
\label{fig:qnnloss}
\end{figure}

\paragraph{The architecture.} In Figure~\ref{fig:qnnarch} the level enters through a cosine embedding, $\cos(i\pi u)$ for $i = 1, \ldots, K$ with $K = 32$, followed by one layer, and the summary passes through two layers. The two embeddings, each of width 64, are multiplied coordinate by coordinate, so that the shape of the posterior, not only its location, can change with $s$, and a two layer head returns the quantile. The two examples of this section train on the check loss alone, with inputs and outputs on their own scales. The worked example of Section~\ref{comp:sec:core} also uses the check loss alone but standardizes inputs and targets by their training means and standard deviations. The horseshoe networks of Section~\ref{comp:sec:horseshoe} use the defaults of the \texttt{gbc} package (Polson and Sokolov 2023), which add an L1 location anchor and a soft penalty on crossing quantiles to the check loss and normalize inputs and outputs.

\begin{figure}[ht]
\centering
\resizebox{\linewidth}{!}{\begin{tikzpicture}[node distance=5mm and 7mm, box/.style={draw,rounded corners=2pt,minimum height=7mm,inner sep=3pt,font=\small,align=center}, ar/.style={-{Stealth[length=2mm]}}]
\node[box] (tau) {$u \sim U(0,1)$};
\node[box, right=of tau] (cos) {cosine\\ embedding};
\node[box, right=of cos] (lt) {layer};
\node[box, below=9mm of tau] (h) {summary $S(D)$};
\node[box, right=of h] (lh1) {layer};
\node[box, right=of lh1] (lh2) {layer};
\node[circle,draw,inner sep=1pt, right=12mm of $(lt)!0.5!(lh2)$] (prod) {$\odot$};
\node[box, right=7mm of prod] (head) {head};
\node[box, right=7mm of head] (q) {$\theta = f_\phi(u, s)$};
\draw[ar] (tau) -- (cos); \draw[ar] (cos) -- (lt); \draw[ar] (lt) -| (prod);
\draw[ar] (h) -- (lh1); \draw[ar] (lh1) -- (lh2); \draw[ar] (lh2) -| (prod);
\draw[ar] (prod) -- (head); \draw[ar] (head) -- (q);
\end{tikzpicture}}
\caption{The quantile neural network. The quantile level and the data summary are embedded separately, combined by a coordinatewise product $\odot$, and mapped to a quantile of the posterior.}
\label{fig:qnnarch}
\end{figure}

\paragraph{Training and sampling.}  Each pass pairs every simulated summary with four fresh uniform levels, in minibatches of 512 simulations, and takes Adam steps on \eqref{gbc:eq:pinball} at learning rate $0.002$ with a cosine decay, for 60 passes. Averaging the network over a grid of levels gives the posterior mean by \eqref{comp:eq:quantileidentity}. Reweighted averages give the distorted values and risk criteria of Section~\ref{gbc:sec:risk} when the utility increases in $\theta$; otherwise they reweight the sorted utilities of the draws.

\paragraph{Crossing and rearrangement.} Nothing forces $f_\phi(u,s)$ to increase in $u$, so a curve read off a grid of levels can cross. Sorting it costs nothing: by the rearrangement inequality\index{rearrangement inequality|textbf} of Hardy, Littlewood and P\'olya (1952) it never raises the integrated check loss at any outcome, nor, by Chernozhukov, Fern\'andez Val and Galichon (2010), the error in estimating a monotone function. Draws $f_\phi(V,s)$ with $V$ uniform already follow the sorted curve, since the law of $f(V)$ has the increasing rearrangement of $f$ as its quantile function, so quantities computed from draws are unaffected. The curves of Figure~\ref{fig:qnnnormal} and the errors reported below are for the raw network output, so sorting could not raise them. 

\paragraph{Normal learning, learned.} Take $\theta \sim N(0, 1)$ and ten observations $y_i \mid \theta \sim N(\theta, 1)$, with the sufficient summary $S = \bar y$. The exact posterior is normal with mean $10\bar y/11$ and standard deviation $1/\sqrt{11} \approx 0.302$. We trained the network on 40{,}000 prior simulations. Figure~\ref{fig:qnnnormal} compares the learned and exact quantile functions at three values of $\bar y$: the mean absolute error over a grid of 99 levels from 0.01 to 0.99, which approximates the Wasserstein distance between learned and exact posteriors, is between 0.007 and 0.013. At $\bar y = 1.5$, 20{,}000 draws from the network have mean 1.359 and standard deviation 0.300, against exact values of 1.364 and 0.302.

\begin{figure}[ht]
\centering
\includegraphics[width=\linewidth]{qnn_normal.pdf}
\caption{Normal learning with a quantile neural network. Left: learned quantile functions (thin lines) and exact posterior quantile functions (thick grey lines) at $\bar y = -1.5$, $0$ and $1.5$. Centre: network draws at $\bar y = 1.5$ against the exact posterior density. Right: the training loss.}
\label{fig:qnnnormal}
\end{figure}

\paragraph{A randomized trial, learned.} Take a two arm trial with 50 patients per arm, success probabilities $p_0$ and $p_1$ with independent uniform priors, and target the risk difference $\tau = p_1 - p_0$. The summary is the pair of observed success rates. We trained on 60{,}000 simulated trials and evaluated the network at 18 control and 27 treated successes. The exact posterior is a difference of two beta variables. The network gives a posterior median of 0.172 and a 90\% interval of $[0.019, 0.332]$; the exact values are 0.174 and $[0.014, 0.328]$. The posterior probability that treatment helps is 0.967 from the network and 0.963 exactly. Figure~\ref{fig:qnntrial} shows both posteriors and a simulation based calibration over 500 simulated trials with 19 draws each; in the saved run, which used two threads, the rank histogram is consistent with uniformity, with a chi squared statistic of 22.2 on 19 degrees of freedom, a value that changes with the thread count.

\begin{figure}[ht]
\centering
\includegraphics[width=\linewidth]{qnn_trial.pdf}
\caption{A two arm trial. Left: network draws of the risk difference at 18 of 50 and 27 of 50 successes, against the exact posterior. Right: ranks of the true risk difference among 19 network draws, over 500 simulated trials; the line is the count expected under perfect calibration.}
\label{fig:qnntrial}
\end{figure}



% comp 05: the algorithm and what it costs; phase 2 rewrite of the moved m14 and m16 text
\section{The algorithm and what it costs}\label{comp:sec:algorithm}\label{gbc:sec:algorithm}

Everything the network knows about the model comes from simulations: the algorithm\index{generative Bayesian computation!algorithm|textbf} is a table of simulated pairs and one regression on it.

\begin{enumerate}[leftmargin=*]
\item \emph{Simulate.} For $j=1,\dots,J$, draw $\theta^{(j)}$ from the prior and a data set $D^{(j)}$ from the model at $\theta^{(j)}$.
\item \emph{Summarize.} Compute $s^{(j)}=S(D^{(j)})$ for every simulated data set, with the code later applied to the observed data.
\item \emph{Train.} Draw levels $u_j$ uniformly on $(0,1)$ and fit $f_\phi$ by minimizing the check loss \eqref{gbc:eq:pinball} over the triples $(\theta^{(j)},u_j,s^{(j)})$, coordinate by coordinate in triangular order if $\theta$ is a vector (Section~\ref{gbc:sec:qnn}).
\item \emph{Infer.} Compute $s=S(D^{\mathrm{obs}})$, draw $U_1,\dots,U_R$ uniformly and return $\theta^{(r)}=f_\phi(U_r,s)$. For a vector $\theta$, draw an independent uniform for each coordinate and evaluate the maps in triangular order, each given the coordinates already drawn.
\end{enumerate}

The $R$ draws approximate the posterior of $\theta$ given $S=s$ (Section~\ref{comp:sec:summaries}); they follow it exactly only if $f_\phi$ is the conditional quantile function, and Section~\ref{comp:sec:checks} says how to test how close it is. 

Training costs $C_{\mathrm{train}}$ for steps 1 to 3; step 4 costs $C_{\mathrm{net}}$ per data set, including $R$ forward passes. A sampler costs $C_{\mathrm{mcmc}}$ per data set. The network is cheaper for $n_{\mathrm{sets}}$ data sets when
\[
C_{\mathrm{train}} + n_{\mathrm{sets}}\,C_{\mathrm{net}} < n_{\mathrm{sets}}\,C_{\mathrm{mcmc}}
\quad\Longleftrightarrow\quad
n_{\mathrm{sets}} > \frac{C_{\mathrm{train}}}{C_{\mathrm{mcmc}} - C_{\mathrm{net}}},
\]
which is close to $C_{\mathrm{train}}/C_{\mathrm{mcmc}}$ when a forward pass costs far less than a chain.

The ten networks of Section~\ref{comp:sec:horseshoe} each took more than an hour to train, so $C_{\mathrm{train}}>10\times3{,}600=36{,}000$ seconds. Break-even is about 57{,}000 data sets at 0.63 seconds per horseshoe fit, 116{,}000 at 0.31 seconds (Section~\ref{comp:sec:cost}) and 440{,}000 at 0.081 seconds (40.6 seconds for 500 data sets, Section~\ref{comp:sec:horseshoe}), or more. The timings come from different machines, so these are orders of magnitude: the sampler wins unless there are at least tens of thousands of data sets.



For a new $s$, rejection ABC\index{approximate Bayesian computation|textbf}\index{ABC|see{approximate Bayesian computation}} keeps the fraction $a$ of the stored summaries nearest to $s$ and reads the posterior off their $\theta^{(j)}$. That is a local, nearest neighbour estimate of the conditional quantile function at $s$ that needs a tolerance and discards $(1-a)J$ simulations at every data set. The network is one global fit to all $J$ pairs. For rejection ABC with the tolerance balanced against Monte Carlo variance, the expected quadratic loss of the posterior mean exceeds its limit by a term of order $J^{-4/(4+k)}$ in the summary dimension $k$ (Fearnhead and Prangle 2012).  No theorem covers a network trained at random levels by stochastic gradient descent, so we claim no advantage in rate.

A network is trusted only where it has seen summaries: outside the simulated region, or where simulations are sparse, it extrapolates and the training loss does not flag it (Section~\ref{comp:sec:checks}).



% comp 06: summaries; phase 2 rewrite (the opening justification sentence of the old SPLIT paragraph is dropped; comp/s04 now says what the proposition gives)
\section{What the network sees: summaries}\label{comp:sec:summaries}

The network never sees the data. It sees $S(D)$, a vector we chose, and learns $\pi(\theta\mid S=s)$ (Proposition~\ref{gbc:prop:pinball}). A summary\index{summary statistics|textbf} is a promise about what we are willing to forget. If $S$ is sufficient the promise costs nothing; if not, this valid posterior conditions on less than the data and is less concentrated on average over data sets. Summarizing is the first of two approximations to $\pi(\theta\mid D^{\mathrm{obs}})$; the second is finite training.

Sufficient summaries lose nothing, as in the normal mean and the two arm trial of Section~\ref{gbc:sec:qnnexplained}, but when the support does not depend on $\theta$, a sufficient statistic of fixed dimension exists, under regularity conditions, only for exponential families. Learned summaries start from Fearnhead and Prangle (2012), who take $S(D)=\E[\theta\mid D]$ and estimate it by regressing simulated $\theta$ on simulated data. By the tower property the posterior given this summary keeps the full posterior mean, and its variance is the average of the full posterior variances over data sets sharing the summary. 

The posterior mean can carry no information. If the prior and the likelihood are unchanged when $\theta$ is replaced by $-\theta$, the posterior is symmetric about zero, typically with a mode on each side, and $\E[\theta\mid D]=0$ for every $D$, so the Fearnhead and Prangle summary is a constant and the network returns the prior. Coverage on held out simulations does not reveal the failure, since intervals from $\pi(\theta\mid S)$ have their nominal coverage over the prior predictive distribution\index{prior predictive distribution|textbf} for every $S$, and the global and local ranks of Section~\ref{comp:sec:checks} are flat for a network that returns the prior. The held out loss against a network given another summary shows it, and so does a reference posterior at the observed summary.

The third kind puts in what a frequentist would compute. For a binary treatment a natural choice collects the weighted outcome means, the covariate sums by arm divided by $N$ (the arm means times the arm shares) and the covariate second moments,
\[
\begin{split}
S(D) = \Bigl(&\tfrac1N\textstyle\sum_i \tfrac{W_iY_i}{\hat e(X_i)},\ \tfrac1N\sum_i\tfrac{(1-W_i)Y_i}{1-\hat e(X_i)},\\
&\tfrac1N\textstyle\sum_i X_iW_i,\ \tfrac1N\sum_i X_i(1-W_i),\ \tfrac1N\sum_i X_iX_i^\top\Bigr),
\end{split}
\]
to which one adds outcome regression coefficients by arm, treated and control means within propensity strata and any consistent estimator of the target. The first two components are inverse probability weighted outcome means (Chapter~\ref{ch:ipw}), whose difference is the Horvitz and Thompson estimate of the average effect. The covariate sums record imbalance; the second moments let the network learn its importance for the outcome model.

Fit $\hat e(X)$ inside $S$ on every simulated data set. If the observed analysis cross fits its nuisances, repeat that cross fit inside every simulated $S$, at the cost of nuisance fits on all $J$ data sets, so the network learns the law of the statistic it will see.

Chapter~\ref{ch:blp} proves what a good summary buys (Proposition~\ref{blp:prop:bvm}). Under regularity conditions the posterior given $S$ is asymptotically normal, centred at the efficient minimum distance estimator built from $S$. If $S$ contains an asymptotically normal estimator of a subvector of $\theta$, the summary posterior for that subvector is asymptotically at least as concentrated as the estimator, with a limiting covariance no larger in the matrix order. A good estimator of the target inside $S$ therefore bounds the loss from summarizing. The conditions include a smooth injective limit of $S$, a prior density continuous and positive at the truth, and conditions on the density of $S$, because convergence in distribution of $S$ does not by itself give convergence of the posterior given $S$. The result concerns the exact summary posterior, not a trained network.

No summary of a covariate can produce heterogeneity in $\tau(x)$ if the outcome model omits its interaction with treatment. Choose the summary with the model.

The second approximation is finite training: the network class, the optimizer and the number of simulations $J$. Padilla, Tansey and Chen (2022) give rates for quantile regression with ReLU networks, for an exact empirical risk minimizer at one fixed quantile level; for targets with compositional structure, an intrinsic dimension replaces the summary dimension $k$. Our network is trained at levels drawn at random, for which we know no rate. Theorem A.4 of Polson and Sokolov (2026a), proved there in outline, gives consistency for fully connected ReLU networks of the level and the summary: for a bounded parameter, summaries in a compact set, a H\"older continuous conditional quantile function and a width that grows slowly enough with $J$, the expected integrated check loss of the empirical minimizer converges to its infimum, the expectation being over the training set as well. It does not cover the product architecture of Section~\ref{comp:sec:network}, whose complexity we have not bounded, and none of it bounds the error of the stochastic gradient descent that trains the network.

% comp s07 (W4, phase 2): the causal core; log misc/work-orders/comp-phase2-W4.applied.md
\section{Posteriors of effects from simulated pairs}\label{gbc:sec:efficiency}\label{comp:sec:core}

The network does not know what an average treatment effect is. It learns the law of whatever the simulator puts in the target, given whatever it puts in the summary, so the causal content of the analysis is decided before training starts, by six decisions\index{generative Bayesian computation!posteriors of effects|textbf}.

\begin{enumerate}[leftmargin=*]
\item \emph{Choose the estimand.} The simulator draws both potential outcomes, so it can compute the model based effect (the average of $\tau(X_i;\theta)$ over the observed covariates), the sample effect (the average of the simulated $Y_i(1)-Y_i(0)$, whose observed halves are part of the simulated data) or the population effect. The posterior of the sample effect depends on the joint law assumed for the pair and need not be wider than that of the model based effect (Section~\ref{sec:imputation}).
\item \emph{Train on it.} Train on $\psi(\theta)$ for the model based or population effect, or the average simulated $Y_i(1)-Y_i(0)$ for the sample effect, to learn its marginal posterior directly. Alternatively, push joint draws of $\theta$ through $\psi$ using the triangular map (Section~\ref{comp:sec:noise}). Separate marginal networks lose dependence and give the wrong variance for functionals of correlated coordinates.
\item \emph{Learn the effect as a function.} Give the network $x$ as an input and draw $x$ at random in the simulator, or train at a grid. A joint statement about the curve needs joint draws of its values at the grid points, from a triangular map over the points (Section~\ref{comp:sec:noise}) or from joint draws of $\theta$; independent levels at each point lose the dependence.
\item \emph{Put the estimator in the summary.} With an AIPW or H\'ajek estimate in $S$ (Chapter~\ref{ch:ipw}) the exact summary posterior is asymptotically at least as concentrated as the estimator (Section~\ref{comp:sec:summaries}). Cross fit its nuisances inside the simulator, since residuals on the training simulations are too small; in Chapter~\ref{ch:blp} cross fitting raised the average coverage of nominal 90 percent intervals from 0.85 to 0.94.
\item \emph{Simulate the confounding we fear.} Calibration on simulated data certifies the computation inside the simulator, not the simulator against the world. The prior predictive must contain confounding of the strength and form the application plausibly has, under the coupled prior of Chapter~\ref{ch:horseshoe}.
\item \emph{Check.} Run the checks of Section~\ref{comp:sec:checks}, against an exact reference where there is one.
\end{enumerate}

\paragraph{What holds.} In a correctly specified regular parametric model, with a prior positive and continuous at the truth and summaries that are sufficient or contain an efficient estimator, the Bernstein von Mises theorem gives a posterior that is asymptotically normal, centred at an efficient estimator, with the inverse Fisher information as variance. By the delta method any smooth functional, such as an average effect $\tau_Q(\theta)$, then has a posterior efficient for that model, and the network inherits this up to its training error, which no theorem bounds.

\paragraph{What does not hold automatically.} Semiparametric efficiency, the bound of Hahn (1998) for the nonparametric model of the regression and the propensity score, attained by the estimator of Hirano, Imbens and Ridder (2003), is another matter. A prior that ignores the propensity score can concentrate around the right regression function and still give an average effect biased at the $\sqrt N$ scale, when the regression is rough and the propensity informative (Ray and van der Vaart 2020), and a shrinkage prior can produce regularization induced confounding (Hahn and coauthors 2018). Generative computation changes how the posterior is computed, not which posterior it is: with such a prior every map returns the same wrong posterior. The repair is in the prior or in the estimator: Chapter~\ref{ch:horseshoe} couples the prior to the treatment model, with the limits that chapter reports for large $p$, and Section~\ref{sec:baipw} puts the horseshoe inside the doubly robust score of DML.

\subsection{A worked example with an exact posterior}

The test model is the stylized supplement trial of Chapter~\ref{ch:cdt}, after the Lanarkshire milk experiment (Section~\ref{sec:milk}): $n=200$ children with a baseline measure of undernourishment $Z_i\sim N(0,1)$ and a supplement $W_i$ given with probability $\expit(1+0.8Z_i)$. Covariates and assignments are drawn once and held fixed: 140 treated and 60 control children, with mean $Z$ of 0.26 and $-0.72$ respectively. The outcome is $Y_i=\mu_w+\alpha_wZ_i+\varepsilon_i$ for $w=W_i$, with $\varepsilon_i\sim N(0,1)$ and the variance known, and $\theta=(\mu_0,\alpha_0,\mu_1,\alpha_1)$ has the prior $N(0,4I)$. The targets are the model based effect $\psi=(\mu_1-\mu_0)+(\alpha_1-\alpha_0)\bar Z$, with $\bar Z$ the mean of the $Z_i$, and the conditional effect $\tau(z)=(\mu_1-\mu_0)+(\alpha_1-\alpha_0)z$.

The four within arm least squares coefficients $S$ are sufficient, so every network error is amortization error (Section~\ref{comp:sec:checks}). The arms have independent normal posteriors: with design $B_w$ (rows $(1,Z_i)$) and least squares coefficients $b_w$,
\[
(\mu_w,\alpha_w)^\top\mid S \sim N\bigl(V_wB_w^\top B_wb_w,\ V_w\bigr),\qquad V_w=\bigl(B_w^\top B_w+\tfrac14 I_2\bigr)^{-1}.
\]
The covariance depends only on the fixed design, so the exact posterior sd of $\psi$, the square root of $(1,\bar Z)(V_0+V_1)(1,\bar Z)^\top$, is 0.179 at every data set, against a prior sd of 2.83.

The ATE network is the quantile network of Section~\ref{comp:sec:network} with $\psi(\theta)$ as the target, trained with the check loss on 100,000 simulated pairs of $\theta$ and $S$ for 40 passes, about 26 seconds on two threads of an Intel server node. The observed data set is one draw of $\theta$ from the prior and the data simulated at it; the checks use 500 fresh studies with 100 draws each and 20,000 fresh pairs. Each network was trained once, so the observed data set results are for one data set and one seed.

At the observed data set the exact posterior of $\psi$ has mean 0.079 and sd 0.179, and 20,000 draws from the network have mean 0.067 and sd 0.185. Their Wasserstein-1 distance is 0.015, or 0.085 of the posterior sd (Figure~\ref{comp:fig:ate}a and c).

\begin{figure}[ht]
\centering
\includegraphics[width=\linewidth,trim=0bp 185bp 0bp 0bp,clip]{gbc_ate.pdf}
\caption{The ATE network of the worked example. (a) Draws of $\psi$ at the observed data set, the exact posterior density and the true value (dashed line). (b) Rank histograms over 500 studies for the ATE network, the exact posterior and the prior only network (trained with the summary replaced by a constant); the line marks the expected count of 50. (c) Quantile functions of the same three at the observed data set. (d) Wasserstein-1 error over the exact posterior sd against the Mahalanobis distance of $S$ from the center of the training summaries, for the 500 studies (the 50 farthest as squares) and along a ray leaving the center in the direction in which the exact mean of $\psi$ moves fastest; the dotted line is the largest distance among the training summaries.}
\label{comp:fig:ate}
\end{figure}

The rank check over 500 studies (Table~\ref{comp:tab:ate}, Figure~\ref{comp:fig:ate}b) cannot tell the ATE network from the exact posterior. Nor can it tell either from a prior only network, which is trained with the summary replaced by a constant and has intervals about 16 times as wide. The held out loss can: it puts the prior only network at 0.798 against 0.050 for the exact quantile function, and resolves the ATE network's excess, 0.00024, at 8.9 standard errors.

\begin{table}[ht]
\centering
\caption{Checks of three posteriors for $\psi$. Ranks: 500 studies with 100 draws each, in ten bins. Intervals run between the 5th and 95th percentiles of the draws (Weibull plotting positions). Held out loss: the check loss on 20,000 fresh pairs with eight random levels each, for the exact quantile function and for the raw network function (0.05027 for the ATE network after sorting).}
\label{comp:tab:ate}
\footnotesize\setlength{\tabcolsep}{4pt}
\begin{tabular}{L{1.2in}>{\raggedleft\arraybackslash}p{0.8in}>{\raggedleft\arraybackslash}p{0.55in}>{\raggedleft\arraybackslash}p{0.55in}>{\raggedleft\arraybackslash}p{0.6in}}
\toprule
 & Rank $\chi^2$ ($p$) & Coverage of 90\% & Width of 90\% & Held out loss \\
\midrule
Exact posterior & 10.12 (0.341) & 0.876 & 0.603 & 0.05012 \\
ATE network & 8.12 (0.522) & 0.870 & 0.601 & 0.05037 \\
Prior only network & 8.36 (0.498) & 0.882 & 9.514 & 0.79757 \\
\bottomrule
\end{tabular}
\end{table}

The 50 studies with summaries farthest from the training center (Mahalanobis distance) have a flat rank histogram and a Wasserstein-1 error of 0.110 of the exact sd, against 0.061 for the other 450, although all 500 lie inside the training range (largest distance 4.59, against 5.61 among the training summaries). Beyond that range, along the ray of Figure~\ref{comp:fig:ate}d, the error passes the posterior sd (1.19 of it at distance 10).

The CATE network, with $z$ as a further input drawn from $N(0,1)$ for each simulated pair, is less accurate than the ATE network at this data set: its Wasserstein-1 errors are 0.114, 0.531 and 0.128 of the exact sd at $z=-1.5$, 0 and 1.5 (Figure~\ref{comp:fig:cate}e).

The 0.531 is mostly one unlucky fit. We trained ten networks of each kind, differing in initial weights and minibatch order, at each of 10{,}000, 40{,}000, 100{,}000 and 400{,}000 simulated pairs (each training set a prefix of the next and the same for every seed, with the architecture and 40 passes above), and compared the sorted network quantiles on 999 levels from 0.001 to 0.999 with the exact posterior quantiles on 2{,}000 fresh studies, by the Wasserstein-1 error over the exact sd. Each network gives a median and a 90th percentile of this error over the studies and a coverage, and we report their means over the ten seeds.

At 100{,}000 pairs the CATE network at $z=0$ has median error 0.15 (0.14 at $z=-1.5$, 0.11 at $z=1.5$), 90th percentile 0.33 and coverage of 90 percent intervals 0.91. At the observed data set its error averages 0.15 over the ten seeds, with standard deviation 0.13 and largest value 0.49, so our single network, trained on a different draw of simulations, lies above all ten. The error falls roughly as $J^{-1/2}$, from a median of 0.60 at 10{,}000 pairs to 0.09 at 400{,}000. The ATE network is more accurate at every size, with median error 0.30 to 0.41 of the CATE network's at $z=0$ (0.054 at 100{,}000 pairs), so a conditional effect needs several times the simulations of an average one for the same accuracy.

\begin{figure}[ht]
\centering
\includegraphics[width=\linewidth,trim=0bp 0bp 0bp 362bp,clip]{gbc_ate.pdf}
\caption{The conditional effect. (e) Exact posterior mean and 90 percent band of $\tau(z)$ at the observed data set, with the median and the 5 and 95 percent points of the CATE network. (f) Posterior sd of $\tau(z)$ from the exact posterior, from the exact marginals combined as if independent, from triangular draws and from independent draws.}
\label{comp:fig:cate}
\end{figure}

The second decision matters for the same curve. The exact posterior correlation of $\Delta\mu=\mu_1-\mu_0$ and $\Delta\alpha=\alpha_1-\alpha_0$ is 0.368. Triangular draws, from a network for $\Delta\mu$ and one for $\Delta\alpha$ given $S$ and the drawn $\Delta\mu$, have correlation 0.390; independent draws take $\Delta\alpha$ from a network given $S$ alone. At $z=-1.5$ and 1.5 the sd of $\tau(z)$ is 0.990 and 1.017 times the exact sd for triangular draws, 1.251 and 0.862 for independent draws, and 1.247 and 0.858 for the exact marginals combined as if independent: the fault is the lost dependence, not the networks (Figure~\ref{comp:fig:cate}f). We chose the assignment intercept so that the components are correlated: at intercept 0 the correlation is 0.001.

The example is narrower than the recipe: the summary is sufficient, the prior normal, the only confounder one observed covariate, and a network serves one design. It has no AIPW summary and no cross fitted nuisances, and we have not tried the coupled horseshoe prior with a network.

% comp section 08
\section{Learning the horseshoe posterior, and why we usually do not}\label{comp:sec:horseshoe}\label{gbc:sec:horseshoe}

The Gibbs sampler of Chapter~\ref{ch:horseshoe} fits a horseshoe regression in 0.31 seconds at $n = 200$ and $p = 51$ (Section~\ref{comp:sec:cost}), so a network has to earn its training time\index{generative Bayesian computation!horseshoe posterior|textbf}. In our book \emph{Generative Bayesian Computation} (Polson and Sokolov 2026a) we gave it a deliberately easy case: ten coefficients, a fixed design of $n = 50$ rows of independent standard normals, known error standard deviation $\sigma = 0.5$, and 80,000 draws from the prior with the data they imply. The least squares estimate $b$ is sufficient for the coefficient vector $\beta$ given the design and $\sigma$, so each network maps $b$ to the quantile function of one $\beta_j$. In this experiment $\beta_j \mid \lambda_j, \gscale \sim N(0, \lambda_j^2\gscale^2)$ with standard half Cauchy local and global scales, not multiplied by $\sigma$ as in Chapter~\ref{ch:horseshoe}. Each network has width 256 and was trained for 3,000 full batch epochs with the defaults of the \texttt{gbc} package, whose loss adds an L1 location anchor and a crossing penalty to the check loss; the reference Gibbs sampler kept 20,000 draws after discarding 2,000.

The first obstacle is scale. The horseshoe prior has tails at least as heavy as a Cauchy's, so under it $\beta_j$ has no mean and no variance. The sample standard deviations of the 80,000 draws run from 808 to 21,740, while the median of $|\beta_j|$ is 0.56. Standardizing by the training mean and standard deviation squashes 92 percent of the draws to within 0.01 of the mean on the standardized scale, and an error of 0.001 on the standardized scale is an error of at least 0.81 in $\beta_j$, more than the median of $|\beta_j|$.

Three pipelines differ only in scale. The raw pipeline maps $b$ to $\beta_j$. The asinh pipeline maps $\operatorname{asinh}(b)$ to $\operatorname{asinh}(\beta_j)$ and returns the quantiles of $\beta_j$ through $\sinh$. The residual pipeline maps $\operatorname{asinh}(b)$ to $\beta_j - b_j$ and adds $b_j$ back to the fitted quantiles. Only the residual target is easy to scale: given $\beta$, the error $b - \beta$ is normal with covariance $\sigma^2(X^\top X)^{-1}$, which does not depend on $\beta$, so $\beta_j - b_j$ is exactly normal under the prior predictive distribution, with standard deviations from 0.067 to 0.096 across the coefficients.

We fixed two criteria before the runs. A pipeline is calibrated if the central 90 percent interval of each of its ten networks contains $\beta_j$ in 85 to 95 percent of 2,000 new data sets. It agrees with the Gibbs sampler of Makalic and Schmidt (2016), whose own intervals, pooled over ten coefficients and 500 data sets, covered 89.9 percent, if for every coefficient the median over 500 data sets of the Kolmogorov and Smirnov distance between the two posteriors is below 0.05. No pipeline passed. Two raw networks returned intervals of zero width that never contained $\beta_j$, and the other eight were 17 to 88 times as wide as the sampler's. The asinh and residual pipelines were calibrated and still disagreed with the sampler, by median distances of 0.16 to 0.18 and 0.062 to 0.078. Coverage averaged over data sets can be right while the posterior for a given data set is wrong. The residual intervals came within 5 percent of the sampler's widths, but their posterior medians were only 1 to 7 percent closer to the truth than the least squares estimates.

The sampler was also far cheaper: 40.6 seconds on one core for the 500 data sets, run 100 at a time by a vectorized sampler, or 0.081 seconds each, while one network took more than an hour to train, so the ten networks of one pipeline cost what the sampler spends on more than 440,000 data sets. 

Training a coupled horseshoe network would mean drawing the local and global scales from half Cauchy distributions, shared between the treatment and outcome equations when the coupling is through scales; leaving $\tau$ and the coefficient of the propensity covariate unshrunk, with wide normal priors; and holding the design at the observed covariates, so that one network serves one study. With $p$ larger than $n$ the least squares estimate is not unique, but for a fixed design $X^\top Y$ is still sufficient for the coefficients given $\sigma$ (with $Y^\top Y$ added when $\sigma$ is unknown); it has $p$ coordinates, and reducing it brings back the choices of Section~\ref{comp:sec:summaries}. We have not tried a network on the coupled design of Chapters~\ref{ch:dml} and~\ref{ch:horseshoe} ($n = 200$, $p = 51$, five confounders).

% comp s09 (W6, phase 2); log misc/work-orders/comp-phase2-W6.applied.md
\section{Checks and failure modes}\label{comp:sec:checks}

A failing sampler leaves traces in its output; a trained network returns a smooth quantile function whatever it has learned, so it needs outside checks\index{generative Bayesian computation!checks and failure modes|textbf}.

Split $\hat R$\index{split R hat@split $\hat R$} and the effective sample size say whether a chain has settled, and are read for the effect, not only for the nuisance parameters. The posterior predictive check\index{posterior predictive check|textbf} asks whether the model could have produced the data on the summaries the effect depends on: outcome means by arm within propensity strata, residual variance by arm, and the tails.

For split $\hat R$, cut each chain in half. With $m$ draws per half, let $W$ be the mean variance within halves and $B$ be $m$ times the variance of half chain means. The posterior variance estimate and diagnostic are
\[
\widehat{\mathrm{var}}^{+}=\frac{(m-1)W+B}{m},\qquad \hat R=\sqrt{\frac{\widehat{\mathrm{var}}^{+}}{W}}.
\]
Values near 1 say the halves agree. The effective sample size is the number of draws divided by $1+2\sum_{k\ge1}\rho_k$, where $\rho_k$ is the autocorrelation at lag $k$ and the sum is truncated as in Geyer's initial positive sequence. The Monte Carlo error of a posterior mean is about the posterior standard deviation over the square root of the effective sample size. For a probability of benefit, apply the same to $h=\one(\tau>0)$, which can mix differently from $\tau$, so a fixed number of sweeps does not fix the precision of a tail quantity.

Vehtari and coauthors (2021) replace the classical $\hat R$ by the larger of a rank normalized split $\hat R$ and its folded version, flag values above 1.01, and add a tail effective sample size from the indicators at the 5th and 95th percentiles; the checks of Section~\ref{sec:hsgibbs} used the classical versions. No diagnostic proves that every mode was visited.

Simulated pairs also test a network, in four checks (Figure~\ref{comp:fig:checks}).

\begin{figure}[ht]
\centering
\resizebox{\linewidth}{!}{\begin{tikzpicture}[node distance=7mm,box/.style={draw,rounded corners=2pt,minimum height=7mm,text width=2.3cm,inner sep=3pt,font=\small,align=center},ar/.style={-{Stealth[length=2mm]}}]
\node[box](a){held out loss};
\node[box,right=of a](b){global ranks};
\node[box,right=of b](c){local ranks at $s$};
\node[box,right=of c](d){reference};
\draw[ar](a)--(b);\draw[ar](b)--(c);\draw[ar](c)--(d);
\end{tikzpicture}}
\caption{The four checks of a trained network, in the order we run them.}
\label{comp:fig:checks}
\end{figure}

\paragraph{Held out check loss.} The first check meters the amortization\index{amortization|textbf} error, the distance at a summary $s$ between the network's posterior and the exact one. For a network that increases in $u$ (sorting, which never raises the loss, makes it so: Exercise~\ref{ex:chcomp:4}), averaging over the prior predictive\index{amortization} gives
\begin{equation}\label{comp:eq:amortization}
\E\int \{F_\phi(t \mid S) - F(t \mid S)\}^2\,dt = 2\{\bar{\mathcal L}(\phi) - \bar{\mathcal L}^\ast\}.
\end{equation}
Here $F_\phi(\cdot\mid s)$ is the distribution function of the draws $f_\phi(U,s)$, $\bar{\mathcal L}(\phi)=\E\,\rho_U\{\theta-f_\phi(U,S)\}$ is the expected check loss and $\bar{\mathcal L}^\ast$ its minimum, reached by the exact posterior. Held out pairs estimate $\bar{\mathcal L}(\phi)$ and so rank networks without a reference posterior, though an average can hide a large error at one summary.

\paragraph{Rank histograms.} The second check, the global ranks, ranks the true $\theta$ among $L$ draws from the network's posterior at $S(D)$: simulation based calibration\index{simulation based calibration|textbf}.

\begin{proposition}[Rank uniformity]\label{comp:prop:ranks}
Let $(\theta, D)$ come from the prior and the model\index{simulation based calibration}, with $\theta$ real valued, and let $S$ be any function of $D$. If the law of $\theta$ given $S$ has no atoms and, given $S$, the $L$ draws are independent of $\theta$ and of each other with that law, then the number $B$ of draws below $\theta$\index{rank statistic|textbf} is uniform on $\{0, \dots, L\}$.
\end{proposition}
\begin{proof}
Given $S$, $\theta$ and the $L$ draws are $L + 1$ independent draws from one law, hence exchangeable and almost surely distinct, so each of the $L + 1$ places $\theta$ can take in their order is equally likely: $\Prob(B = b \mid S) = 1/(L + 1)$ for every $b$.
\end{proof}

Even a network returning the prior has flat ranks and nominal held out coverage (Exercise~\ref{ex:chcomp:1}). The check certifies calibration\index{calibration|textbf} over the prior predictive, not sharpness or accuracy at the observed data. U shaped histograms mean draws too narrow; humped, too wide; tilted, shifted.

\paragraph{Local ranks.} The third check repeats the test at the observed summary $s$, for $M$ simulations and $M$ evaluations of the network: draw $M$ parameters from its posterior at $s$, simulate a data set from each, recompute the posterior, and rank the generating value. The test is not the same, however, because the parameters come from the posterior at $s$ and not from the prior: Proposition~\ref{comp:prop:ranks} does not apply, the local ranks of the exact posterior are not flat in general (in the normal example of Section~\ref{comp:sec:network} at $\bar y = 1.5$ about two thirds of them fall in the upper half), and those of a network that returns the prior are exactly flat at every $s$ (Exercise~\ref{ex:chcomp:1}). The check is informative when the prior is weak relative to the data, and with an informative prior it can pass a useless network and fail the exact posterior. In the trial, each draw of the risk difference $\tau$ is paired with $p_0$ from its conditional posterior given $\tau$ and the data, and $p_1 = p_0 + \tau$, with no clipping; a draw of $\tau$ outside $(-1, 1)$ has no valid pair and is redrawn. Importance weights proportional to likelihood times prior over network density target the exact posterior, provided these quantities are evaluable and the network density is positive on the posterior support; we have not tried this correction. 

\paragraph{Evidence.} We ran the local check with $M = 400$ and 19 draws each, at two threads, for a network retrained with the same design, at the observed trial of Section~\ref{gbc:sec:qnnexplained}, 18 of 50 and 27 of 50 successes, and at an extreme trial, 1 of 50 and 49 of 50. Against the exact posterior's expected counts under the same protocol, the network's counts give chi squared 21.1 on 19 degrees of freedom ($p = 0.33$) at the typical trial and 28.4 ($p = 0.08$) at the extreme one, so neither rejects agreement with the exact reference.

Against a flat histogram they give 24.0 ($p = 0.20$) and 222 ($p < 10^{-35}$). The extreme histogram piles up at the top ranks, but at so extreme a summary the exact posterior's own local ranks do so too, for the reason given above: with the exact posterior in place of the network, 200 repeats of the protocol give a chi squared statistic with median 20 at the typical trial and 244 at the extreme one. A direct comparison of the posteriors finds what the ranks miss. At the typical trial the network's posterior is within 0.003 of the exact one in Wasserstein distance; at the extreme trial it is shifted down and too wide, its 5\% quantile 0.834 against the exact 0.853 and its Wasserstein distance 0.009, because few training simulations land near a risk difference of 0.9 (Figure~\ref{fig:openlocalsbc})\index{simulation based calibration!local}. The global ranks of the network of Section~\ref{gbc:sec:qnnexplained} passed (chi squared 22.2 at two threads, Figure~\ref{fig:qnntrial}), since a risk difference above 0.9 has prior probability 0.005. Problem 6 of Chapter~\ref{ch:open} asks for a bound on such local error.

\begin{figure}[ht]
\centering
\includegraphics[width=\linewidth]{open_localsbc.pdf}
\caption{Local simulation based calibration of the trial posterior at a typical and an extreme data set: rank histograms of the network ($M = 400$ simulations, 19 draws each, network trained on 60{,}000 trials) per 400 draws, with the rank histogram of the exact posterior under the same protocol (black steps) and the flat count (dotted line).}
\label{fig:openlocalsbc}
\end{figure}

\paragraph{Reference.} Compare with a trusted posterior: the trial's beta difference, the normal posterior (Section~\ref{comp:sec:core}) or the horseshoe Gibbs sampler (Section~\ref{sec:hsgibbs}). The distance is local amortization error; the reference loss gives $\bar{\mathcal L}^\ast$. A tractable cousin tests the architecture, not the application.

Known failures, mostly from the experiments in our book \emph{Generative Bayesian Computation} (Polson and Sokolov 2026a):
\begin{itemize}
\item \emph{Under-dispersion.} On three of the four tasks we took from the benchmark of Lueckmann and coauthors (2021), separate networks trained with the check loss alone on 10,000 simulations gave central 90 percent marginal intervals that contained 31 to 52 percent of the true parameter components over the ten benchmark observations of each task (31 and 33 of 100, 26 of 50), for reasons not diagnosed.
\item \emph{High signal to noise.} Near deterministic outcomes, such as the luminosities of Chapter~\ref{ch:gente}, are the risk. The networks of that book gave intervals too wide in simulations and too narrow on its lake data set; for the narrow ones it offers only a heuristic, that the product of level and summary features starves the spread. Remedies: heteroscedastic ensembles, conformal calibration (marginal coverage only).
\item \emph{Lost dependence.} Separate networks give the marginals without the dependence, so a contrast gets the wrong variance (Exercise~\ref{ex:chcomp:2}), which only its own ranks reveal; the triangular map of Section~\ref{comp:sec:noise} keeps it.
\item \emph{Blurred jumps.} In surrogate regression of a response surface with a jump, averaging an ensemble blurs the jump, since its members place it at different points; one network given an estimated side of the jump did better. We have not tested a jump in a posterior or an outcome surface at a regression discontinuity.
\end{itemize}

% comp s10 (W6, phase 2); log misc/work-orders/comp-phase2-W6.applied.md
\section{Choosing the engine}\label{comp:sec:engine}

Table~\ref{tab:cost} compares engines and cost orders for $n$ units and $p$ covariates\index{engine choice|textbf}; constants depend on the machine.

\begin{table}[ht]
\centering
\caption{Engines by situation. $N_{\mathrm{it}}$ is the number of iterations and $J$ the number of simulated pairs.}
\label{tab:cost}
\footnotesize\setlength{\tabcolsep}{3pt}
\begin{tabular}{L{0.8in}L{0.62in}L{0.82in}L{1.0in}L{0.86in}}
\toprule
Situation & Engine & What it needs & How it fails & Cost order \\
\midrule
Gaussian outcome & Horseshoe Gibbs & Closed form full conditionals & Scales mix slowly & $np^2 + N_{\mathrm{it}}p^3$; $N_{\mathrm{it}}n^2p$ if $p > n$ \\
Binary or count outcome & Polya Gamma Gibbs & Logistic or count likelihood & As above & $N_{\mathrm{it}}(p^3 + np^2)$ \\
No closed form conditionals & HMC & Log posterior and its gradient & Funnels (divergences); heavy tails & $N_{\mathrm{it}}\times$ leapfrog steps $\times\, np$ \\
Near normal posterior & Laplace & Likelihood, gradients & Skewed or bounded posteriors, tails & One optimization and one Hessian \\
Speed, very large $n$ & Variational & Likelihood, a tractable family & Tails; error not reported & One optimization \\
Factorized family & Mean field & Independent blocks & Intervals too short & As variational \\
Likelihood unavailable or costly & Trained network & A simulator; no likelihood or gradients & Extrapolation, under-dispersion, lost dependence & $J$ simulations and a fit, then a forward pass per draw \\
One effect, no posterior & DML, $K$ folds & Cross fitted regressions & Interval only; needs overlap, nuisance rates & $2K$ cross validated lasso fits \\
\bottomrule
\end{tabular}
\end{table}

For Gaussian or Polya Gamma conditionals the Gibbs sampler is the default, and for the horseshoe Section~\ref{comp:sec:horseshoe} reports that trained networks lost to it. A network earns its training where the likelihood cannot be evaluated or is expensive, as in the merger studies of Chapter~\ref{ch:blp}, and over many data sets from one design\index{generative Bayesian computation!cost} (Section~\ref{comp:sec:algorithm} gives the break-even).

\section{In practice}

\paragraph{The target.} Every target in this book is a function of the draws $\theta^{(s)}$: an average effect is $\tau_Q(\theta^{(s)})$ for each draw, a merger threshold or a probability of causation another function of the same draws. The engine decides how fast the draws arrive and how closely they follow the posterior; the model, the prior and the summary decide which posterior that is.

\paragraph{What the analyst does differently.} For the horseshoe model she uses the fast sampler when $p$ exceeds $n$, and includes an unpenalized intercept or centres the outcome and every column. She runs 1,500 to 3,000 iterations, drops the first third, and reports split $\hat R$, the effective sample size for the effect and a posterior predictive check. For a network she reports the summary, the held out loss, the ranks of the effect, the local ranks and a reference where one exists.

\paragraph{Against DML.} DML is cheaper in principle, but in the two Gaussian outcome designs, with 51 and 56 columns and 200 and 9,915 units, we timed the horseshoe sampler at half the time of a cross fitted lasso and a twentieth of the time of a cross fitted forest (Section~\ref{comp:sec:cost}). As $p$ grows the factorization dominates and the comparison turns. In the replications of Section~\ref{sec:ricsim}, with $n = 200$, one coupled horseshoe fit, its lasso fit of the treatment included, took 1.6 seconds on average at $p = 150$ and 5.7 seconds at $p = 300$, while everything in a replication except its three horseshoe fits and the six lasso fits of the treatment that feed two of them, DML's ten lasso fits included, took 3.2 and 2.3 seconds, an upper bound on the cost of DML in that design. The reduced form sampler of Section~\ref{sec:crit:bayes} took 9.8 seconds for one chain of 2,000 sweeps at $n = 300$ and $p = 400$, and about 200 seconds for the four chains run to convergence that Table~\ref{tab:crit:head} uses, against 1.0 for a cross fitted lasso. We did not time the Polya Gamma sampler for a binary outcome, whose cost per iteration grows with $n$. DML with a lasso nuisance keeps a cost advantage at large $p$, and may keep one for a binary outcome and in a very large analysis with a nonlinear nuisance that a Gaussian or Polya Gamma model does not fit. In every design the posterior comes with checks for which DML has no counterpart. A trained network repays its training only over many data sets from one design.


\section*{Notes}
Noise outsourcing is a case of the transfer theorem of Kallenberg (2002); Villani (2009) treats optimal transport. Chernozhukov, Fern\'andez Val and Galichon (2010) show that sorting an estimated monotone function never raises its error in an $L^p$ norm. Learning a posterior as a quantile function of noise is the generative Bayesian computation of Polson and Sokolov (2023). Distributional reinforcement learning begins with Bellemare, Dabney and Munos (2017), and quantile regression for it is in Dabney, Rowland, Bellemare and Munos (2018).

Jiang and coauthors (2017) learned summaries with a neural network. Rank checks go back to Cook, Gelman and Rubin (2006) and were developed as simulation based calibration by Talts and coauthors (2018); Lueckmann and coauthors (2021) is the benchmark behind the under-dispersion result. That result, the other failure modes of Section~\ref{comp:sec:checks} and the consistency statement of Section~\ref{comp:sec:summaries} come from our book \emph{Generative Bayesian Computation} (Polson and Sokolov 2026a). The fast sampler for $p > n$ is Bhattacharya, Chakraborty and Mallick (2016), the simple horseshoe sampler Makalic and Schmidt (2016), and the Polya Gamma augmentation Polson, Scott and Windle (2013).

\begin{exercises}
\item\label{ex:ch14:1} Derive the Bhattacharya, Chakraborty and Mallick sampler for $N(A^{-1}Z^\top Y, \sigma^2 A^{-1})$, with $A = Z^\top Z + \Lambda^{-1}$ and $\Lambda$ diagonal as in Section~\ref{comp:sec:cost}, using the Woodbury identity.
\item\label{ex:ch14:2} Verify the Polya Gamma identity of Section~\ref{comp:sec:cost} for $b = 1$ by computing the Laplace transform of the Polya Gamma law defined in Appendix~\ref{app:toolkit}.
\item\label{ex:ch14:3} Show that, for a scalar parameter $\theta$ whose log posterior given $n$ observations $D$ is smooth with a unique interior mode, the Laplace approximation to the normalizing constant $\int p(D \mid \theta)\,\pi(\theta)\,d\theta$ has relative error $O(n^{-1})$ under regularity conditions. Show with a Gamma posterior that the normal approximation to the posterior density has relative error of order $n^{-1/2}$ at a fixed number of standard deviations from the mode.
\item\label{ex:ch14:4} Design summary statistics for an amortized posterior of the average treatment effect in a partially linear model, and argue for their approximate sufficiency.
\item\label{ex:ch14:5} Implement simulation based calibration for a proper version of the horseshoe regression of Section~\ref{sec:hsgibbs} with $n = 100$ and $p = 20$: a fixed design whose first column is a binary treatment and whose other columns are independent standard normals, the prior $\tau \mid \sigma^2 \sim N(0, 100\sigma^2)$ on the treatment coefficient in place of the flat one, the horseshoe on the other coefficients and an inverse gamma prior with shape 2 and scale 1 on $\sigma^2$, all used in the full conditionals. In each of 1,000 replications draw the parameters from this prior, simulate $Y$, run the sampler, keep 19 draws spaced far enough apart to be nearly independent for both $\tau$ and $\gscale$, and rank the generating values. Report the two rank histograms and how you checked the dependence of the kept draws.
\item\label{ex:chgbc:3} Let $Y_1, \ldots, Y_N$ be independent $N(\theta, \sigma^2)$ given $\theta$, with $\sigma$ known and $\theta \sim N(\mu_0, \sigma_0^2)$, as in the normal learning example of Section~\ref{gbc:sec:normal}. Given the summary $s = \bar Y$ the posterior of $\theta$ is $N(\mu_N(s), \sigma_N^2)$ with $\mu_N(s) = (\mu_0/\sigma_0^2 + Ns/\sigma^2)/(1/\sigma_0^2 + N/\sigma^2)$ and $\sigma_N^2 = (1/\sigma_0^2 + N/\sigma^2)^{-1}$. Verify that the check loss at level $u$ is minimized by $\mu_N(s) + \sigma_N\Phi^{-1}(u)$. Then, for the report $a^\ast = \mu_N(s)$ under the loss $(a^\ast-\theta)^2$, compute the mean of the upper 10\% of the loss over the posterior of $\theta$, which is the negative of the conditional value at risk (CVaR) at level 0.1 of the utility $-(a^\ast-\theta)^2$, the mean of the worst 10\% of that utility.
\item\label{ex:ch08:1} Prove Theorem~\ref{comp:thm:outsourcing} for real valued $Y$ by taking $G(u,x)$ to be the conditional quantile function, including laws with atoms.
\item\label{ex:ch08:2} Show that minimizing the expected check loss $\E\rho_u(Y - q)$ over $q$ gives the $u$ quantile of $Y$.
% comp exercises 1 to 4 (W6, phase 2), input by mcomp.tex inside the exercises environment
\item\label{ex:chcomp:1} In the normal example of Section~\ref{comp:sec:network} ($\theta \sim N(0, 1)$, ten observations $N(\theta, 1)$, summary $\bar y$), take a network that returns the prior $N(0, 1)$ at every summary. Show that its rank histogram is flat, over the prior predictive and at any observed summary. Then compute its expected check loss at level $u = 1/2$ and that of the exact posterior, and say which of the four checks of Section~\ref{comp:sec:checks} tell the two networks apart.
\item\label{ex:chcomp:2} At a covariate value $x$ the posterior draws of the two potential outcome means $\mu_1(x)$ and $\mu_0(x)$ have variances $v_1$ and $v_0$ and covariance $c$, which is not zero when the arms share parameters, as under a common prior scale. Show that independent draws of the two, as from separate networks, give $\mathrm{Var}\,\tau(x) = v_1 + v_0$ for $\tau(x) = \mu_1(x) - \mu_0(x)$, while joint draws give $v_1 + v_0 - 2c$. For which sign of $c$ do separate networks overstate the uncertainty about the effect, and how does the triangular map of Section~\ref{comp:sec:noise} repair it?
\item\label{ex:chcomp:3} Use the notation of Section~\ref{comp:sec:algorithm}. (a) The Gibbs sampler of Section~\ref{comp:sec:horseshoe} took 0.081 seconds per data set and training took at least $36{,}000$ seconds; taking $C_{\mathrm{net}}$ as negligible, how many data sets repay the training? (b) The local ranks of Section~\ref{comp:sec:checks} are run on every data set the network analyses and cost $C_{\mathrm{chk}}$ each. Show that the break-even becomes $C_{\mathrm{train}}/(C_{\mathrm{mcmc}} - C_{\mathrm{net}} - C_{\mathrm{chk}})$, and find the largest $C_{\mathrm{chk}}$ for which the network ever repays its training.
\item\label{ex:chcomp:4} Let $q$ be a measurable, integrable function on $(0, 1)$, $q^\ast$ its non-decreasing rearrangement (the quantile function of $q(V)$ for $V$ uniform) and $y$ a real number. Show that $\int_0^1 \rho_u\{y - q^\ast(u)\}\,du \le \int_0^1 \rho_u\{y - q(u)\}\,du$, using $\rho_u(y - q) = (q - y)_+ - u(q - y)$, the fact that $\int_0^1 (q(u) - y)_+\,du$ depends on $q$ only through the law of $q(V)$, and the rearrangement inequality of Hardy, Littlewood and P\'olya (1952), $\int_0^1 u\,q(u)\,du \le \int_0^1 u\,q^\ast(u)\,du$. Conclude that sorting the output of a network over $u$ never raises its expected check loss.

\end{exercises}
